\documentclass[UTF8,10pt,letterpaper,fontset=fandol]{ctexart}
\usepackage[margin=1in]{geometry}
\usepackage{amsmath,amssymb,enumitem,needspace}
\usepackage[unicode,hidelinks]{hyperref}
\pagestyle{plain}
\raggedbottom
\setlength{\parindent}{0pt}
\setlength{\parskip}{1pt}
\setlength{\emergencystretch}{2em}
\linespread{1.1}
\AtBeginDocument{%
\setlength{\abovedisplayskip}{6pt}%
\setlength{\belowdisplayskip}{6pt}%
\setlength{\abovedisplayshortskip}{3pt}%
\setlength{\belowdisplayshortskip}{3pt}}
\setcounter{secnumdepth}{0}
\setlist[enumerate]{label=(\arabic*),leftmargin=2em,itemsep=0.25em,topsep=0.2em}
\ctexset{section={format=\Large\bfseries,beforeskip=1.2ex,afterskip=0.8ex}}
\newcommand{\C}{\mathbb C}
\newcommand{\F}{\mathbb F}
\begin{document}
\begin{center}
{\LARGE 作业 7}\\[0.6em]
2026 年 10 月 28 日\\[0.25em]
姓名：YuanHao Tang\quad 学号：3120265621
\end{center}

\section{第 1 题}
设 $\mathcal V$ 是域 $\F$ 上的线性空间，$\mathcal A\in L(\mathcal V)$，$\mathcal W\subset\mathcal V$ 是 $\mathcal A$-不变子空间，$\psi(\lambda)$ 是 $\mathcal A$ 在 $\mathcal W$ 上的最小多项式。证明存在 $x\in\mathcal W$，使 $\mathcal A$ 在 $x$ 处的最小多项式等于 $\psi(\lambda)$。

\par\Needspace{4\baselineskip}\textbf{解答}\par
若 $\mathcal W=\{0\}$，则 $\psi=1$，取 $x=0$。否则写 $\psi=\prod_{i=1}^r p_i^{m_i}$，其中 $p_i$ 为互异的首一不可约多项式，令 $q_i=\psi/p_i^{m_i}$。由 $\psi$ 的最小性，可选 $v_i\in\mathcal W$ 使 $(\psi/p_i)(\mathcal A)v_i\ne0$。置 $y_i=q_i(\mathcal A)v_i\in\mathcal W$，则

\[
p_i(\mathcal A)^{m_i}y_i=0,\qquad p_i(\mathcal A)^{m_i-1}y_i\ne0.
\]
故 $y_i$ 处的最小多项式为 $p_i^{m_i}$。这些多项式两两互素，由互素分量求和的性质，

\[
\boxed{x=\sum_{i=1}^r y_i\in\mathcal W,\qquad m_x=\prod_{i=1}^r p_i^{m_i}=\psi}.
\]

\section{第 2 题}
设 $\mathcal A\in L(\C_n[\lambda])$ 定义为

\[
\mathcal A[f(\lambda)]=f(\lambda+a),\qquad\forall f(\lambda)\in\C_n[\lambda].
\]
其中 $a\in\C$ 为非零常数；它关于基 $1,\lambda,\lambda^2,\ldots,\lambda^n$ 的矩阵 $A$ 已在第五次作业中求得。求：

\begin{enumerate}
\item 特征矩阵 $\lambda E-A$ 的 Smith 标准形 $S$，以及 $A$ 的 Jordan 标准形 $J$；

\item 使 $X^{-1}AX=J$ 的过渡矩阵 $X$。

\end{enumerate}

\par\Needspace{4\baselineskip}\textbf{解答}\par

\begin{enumerate}
\item 把作用对象的变量记作 $s$，令 $\Delta=\mathcal A-I$。$\Delta$ 降低非恒定多项式的次数，且 $\Delta^{n+1}=0$、$\Delta^n(s^n)=n!a^n\ne0$。因此只有一个大小为 $n+1$ 的 Jordan 块：

\[
\boxed{S=\operatorname{diag}(I_n,(\lambda-1)^{n+1}),\qquad J=J_{n+1}(1)}.
\]

\item 取下列多项式，并令 $X$ 的第 $k$ 列为 $f_k$ 关于基 $(1,s,\ldots,s^n)$ 的坐标（行列从 $0$ 编号）：

\[
f_0(s)=1,\qquad
f_k(s)=\frac{1}{k!a^k}\prod_{r=0}^{k-1}(s-ra)\quad(1\le k\le n).
\]

\[
\boxed{X_{ik}=[s^i]f_k(s)\quad(0\le i,k\le n)},\qquad
X_{kk}=\frac{1}{k!a^k}\ne0.
\]
其中 $[s^i]$ 表示取系数。$X$ 上三角且可逆；由 $f_k(s+a)=f_k(s)+f_{k-1}(s)$（约定 $f_{-1}=0$），得 $AX=XJ$，即 $X^{-1}AX=J$。

\end{enumerate}

\newpage

\section{第 3 题}
设 $E$ 为适当阶数的单位矩阵，下列多项式是 $\lambda E-A$ 的初等因子：

\[
\begin{gathered}
\lambda+3,\quad\lambda+3,\quad(\lambda+3)^2,\quad\lambda+2+5j,\quad\lambda+2-5j,\\
(\lambda+2+5j)^2,\quad(\lambda+2-5j)^2.
\end{gathered}
\]


\begin{enumerate}
\item 求 $\lambda E-A$ 的 Smith 标准形 $S$、$A$ 的自然标准形 $N$ 和 Jordan 标准形 $J$。

\item 令 $\lambda_1=-3$、$\lambda_2=-2+5j$、$\lambda_3=-2-5j$，并定义

\[
T(\lambda_i)\triangleq\{x:\exists k\in\mathbb Z_+,\ (A-\lambda_iE)^kx=0\}.
\]
$\mathbb Z_{+,0}$ 表示非负整数集。证明 $T(\lambda_i)=\ker[(A-\lambda_iE)^2]$，并求 $\dim T(\lambda_i)$，$i=1,2,3$。

\end{enumerate}

\par\Needspace{4\baselineskip}\textbf{解答}\par
以下 $j^2=-1$，$J_m(\rho)$ 的对角元为 $\rho$、上次对角线为 $1$；$C(p)$ 是列伴随矩阵（次对角线为 $1$，末列为首一多项式系数按低次到高次取负）。

\begin{enumerate}
\item 令 $x=\lambda+3$、$q=(\lambda+2)^2+25$。按各根的指数排序并左补零，得到非常数不变多项式 $x,xq,x^2q^2$，次数为 $1,3,6$。因此

\[
\boxed{S=\operatorname{diag}(I_7,x,xq,x^2q^2),\qquad N=\operatorname{diag}(-3,C(xq),C(x^2q^2))}.
\]
记 $K_\pm(a,b)=\operatorname{diag}(a\pm bj,J_2(a\pm bj))$，则

\[
\boxed{J=\operatorname{diag}(-3,-3,J_2(-3),K_-(-2,5),K_+(-2,5))}.
\]

\item 各特征值的最大 Jordan 块均为二阶，故在该根的广义特征空间上 $(A-\lambda_iI)^2=0$；在其他根的分量上 $A-\lambda_iI$ 可逆。因此

\[
\boxed{T(\lambda_i)=\ker(A-\lambda_iI)^2,\qquad (\dim T(\lambda_1),\dim T(\lambda_2),\dim T(\lambda_3))=(4,3,3)}.
\]

\end{enumerate}

\newpage

\section{第 4 题}
已知矩阵 $A$ 的初等因子为

\[
\begin{gathered}
(\lambda+1),\quad(\lambda+1)^2,\quad(\lambda+1+j),\quad(\lambda+1-j),\\
(\lambda+1)^2,\quad(\lambda+1+j)^2,\quad(\lambda+1-j)^2.
\end{gathered}
\]

\begin{enumerate}
\item 求 $\lambda E-A$ 的 Smith 标准形 $S$；

\item 求 $A$ 的自然标准形 $N$、复 Jordan 标准形 $J_{\rm complex}$ 和实 Jordan 标准形 $J_{\rm real}$；

\item 特征值记为 $\lambda_1=-1$、$\lambda_2=-1-j$、$\lambda_3=-1+j$，求各自的几何重数 $q_i$ 与代数重数 $\widehat n_i$；

\item 求 $\dim T(\lambda_i)$，其中

\[
T(\lambda_i)\triangleq\{x:\exists k\in\mathbb Z_+,\ (A-\lambda_iE_n)^kx=0\}.
\]

\end{enumerate}

\par\Needspace{4\baselineskip}\textbf{解答}\par
令 $x=\lambda+1$、$q=x^2+1$，沿用第 3 题的块记号。

\begin{enumerate}
\item 非常数不变多项式为 $x,x^2q,x^2q^2$，总维数为 $11$，所以

\[
\boxed{S=\operatorname{diag}(I_8,x,x^2q,x^2q^2)}.
\]

\item
\[
\boxed{N=\operatorname{diag}(-1,C(x^2q),C(x^2q^2))},
\]

\[
\boxed{J_{\rm complex}=\operatorname{diag}(-1,J_2(-1),J_2(-1),K_-(-1,1),K_+(-1,1))}.
\]
令 $R=\begin{pmatrix}-1&1\\-1&-1\end{pmatrix}$、$R_2=\begin{pmatrix}R&I_2\\0&R\end{pmatrix}$，共轭块配对后得到

\[
\boxed{J_{\rm real}=\operatorname{diag}(-1,J_2(-1),J_2(-1),R,R_2)}.
\]

\item 几何重数是该根的块数，代数重数是块长之和：

\[
\begin{array}{c|cc}\lambda_i&q_i&\widehat n_i\\\hline -1&3&5\\-1-j&2&3\\-1+j&2&3\end{array}
\]

\item 广义特征空间的维数等于代数重数，故

\[
\boxed{\dim T(\lambda_1)=5,\qquad\dim T(\lambda_2)=3,\qquad\dim T(\lambda_3)=3}.
\]

\end{enumerate}

\end{document}
